% BANKS-SOURCE: pdf=203 print=193 kind=section id=9.14
\subsection{Renormalization and symmetry}

The physical theory is taken in $D=4$ space-time dimensions.  Euclidean loop
integrals are continued to $d_{\mathrm{reg}}$ dimensions when dimensional
regularization is used.

% BANKS-SOURCE: pdf=203 print=193 kind=subsection id=9.14.1
\subsubsection{You don't need symmetry}

% BANKS-SOURCE: pdf=203 print=193 kind=source-page
% BANKS-SOURCE: pdf=203 print=193 kind=prose
Wilson's approach to the interplay between renormalization and symmetry is
based on the observation that the relevant and marginal perturbations of a
Lagrangian with symmetry are generally finite in number. At tree level, one can
set the coefficients of all symmetry-breaking operators to zero. Even if one
ignores the symmetry in constructing a cut-off version of the theory, one can
still imagine tuning the coefficients of the symmetry-breaking relevant and
marginal operators to restore the naive Ward identities at the level of
renormalized Green functions.

Wilson's tuning procedure ignores the issue of naturalness that we raised
above.\footnote[18]{Wilson was one of the very first advocates of naturalness as a
criterion for a good theory.} More interestingly, we have encountered a class
of examples in which it fails. These are theories with ``anomalies.'' Certain
terms that appear in loop corrections to the classical Ward identities cannot
be canceled out by adding local counterterms to the action. Anomalies are
always associated with particles whose mass is zero, or can be made to go to
zero by tuning the expectation value of a low-energy field. The only exception
to this is classical scale invariance. In zeroth-order perturbation theory, the
only requirement is that all operators in the Lagrangian be marginal. However
we have seen that most marginal operators are actually marginally relevant or
irrelevant. Interacting scale-invariant theories are scarce, and typically
appear only at isolated points in the space of field theories.

% BANKS-SOURCE: pdf=203 print=193 kind=subsection id=9.14.2
\subsubsection{But symmetries are nice}

On the other hand, we have seen that invoking symmetries makes the whole
process of renormalization less arduous. If we can invent a regulator that
preserves a symmetry we should do it, if only because it saves work in
calculation. Quite frankly, it wasn't until 't Hooft and Veltman introduced
dimensional regularization that calculations in non-abelian gauge theory became
sufficiently transparent for one to decide whether the theory was renormalizable.
The Wilsonian approach can work, but only if one is a master calculator, who
doesn't make mistakes.

More importantly, symmetries can help us to understand the smallness or absence
of terms in the effective action that we might alternatively think of as finely
tuned relevant

% BANKS-SOURCE: pdf=204 print=194 kind=source-page
parameters. This is the philosophy of ``technically natural field theory''
which we have adumbrated above. Indeed, it can be argued that much of the power
of effective field theory comes from its exploitation of symmetry. For example,
we are able to calculate a large number of low-energy scattering amplitudes
involving pions in terms of a small number of parameters just by using chiral
symmetry and effective field theory. Effective field theory alone would give us
very little information.

% BANKS-SOURCE: pdf=204 print=194 kind=subsection id=9.14.3
\subsubsection{And sometimes you get them for free}

The deepest remark one can make about the relationship between symmetry and
renormalization is that symmetries can be emergent. Like much else in this
field, this remark is due to Wilson. To make the point in a simple context,
consider a lattice version of a scalar field action, thought of as a
perturbation of the Gaussian fixed point. There are lots of allowed operators
(e.g. $\sum_\mu(\partial_\mu\phi)^4$) in the theory, on a cubic lattice, which
break continuous rotation invariance. However, for most regular lattices, the
lattice symmetries are sufficient to remove all relevant and marginal operators
that break the continuous symmetry. Thus we can view rotation invariance (and,
after Wick rotation, Lorentz invariance) as an accidental low-energy
consequence of the fact that IR physics is dominated by a fixed point. In fact,
most of the fixed points describing second-order transitions in condensed-matter
physics have emergent rotation invariance, despite coming from underlying
systems that have no such symmetry.

It is likely that Lorentz invariance is not accidental, but rather a consequence
of the theory of quantum gravity. It is a large-distance manifestation of the
underlying gauge principle of general covariance. However, there is lots of
room for emergent continuous symmetries to play a role in our theory of the real
world. In particular, the standard model has a host of exact and approximate
global continuous symmetries like baryon number, lepton numbers, and various
quark flavor numbers. There are good arguments that none of these symmetries is
exact, but they might be emergent infrared symmetries of an underlying theory
in which e.g. only discrete subgroups of them are exact. A detailed exploration
of theories with emergent continuous symmetries is beyond the scope of our
discussion here, but it is an important theme in the exploration of physics
beyond the standard model.

Important examples of this kind of emergent symmetry abound in the standard
model. For example, baryon and lepton numbers, and lepton flavor, are preserved
automatically by the most general renormalizable Lagrangian with the
standard-model gauge symmetries.\footnote[19]{Actually, $B$ and $L$ are broken,
and only $B-L$ preserved, by an anomaly in the standard model. However, at low
energies the breaking is exponentially small.} Furthermore, any theory of
baryon-number violation can be parametrized by the coefficients of a small
number of dimension-6 operators.

A more subtle emergent symmetry is the extra SU(2) in the Higgs sector of the
standard model, which we called custodial SU(2). It is broken by the U(1) gauge
interaction,

% BANKS-SOURCE: pdf=205 print=195 kind=source-page
but that is the only renormalizable interaction which can break it. As a
consequence, the tree-level broken-symmetry relation
% BANKS-SOURCE: pdf=205 print=195 kind=display
\[
  M_W=M_Z\cos\theta_W
\]
can get only finite corrections, when it is expressed in terms of the
renormalized couplings.

% BANKS-SOURCE: pdf=205 print=195 kind=subsection id=9.14.4
\subsubsection{Renormalization and spontaneous symmetry breaking}

In the previous section we discussed the relation between renormalization and
explicit breaking of symmetries. What then of theories with spontaneously
broken symmetry? Our discussions of spontaneously broken symmetry and of the
renormalization group show that the two subjects are really decoupled from each
other. Renormalization is the procedure for deriving the effective theory of
long-wavelength degrees of freedom of a system, while spontaneous symmetry
breakdown is a property of the solution of that effective field theory.

In particular, if we consider the system in a large but finite volume, then
there is no spontaneous breaking of symmetry. There is a unique, symmetric
ground state. The renormalization procedure involves integrating out
short-wavelength degrees of freedom, which, because of the approximate locality
of interactions, are insensitive to the volume of the system.

A procedure for seeing this decoupling explicitly in perturbative calculations
is to calculate the quantum action, rather than connected Green functions. One
can then test the system for spontaneous symmetry breaking by looking to see
whether there are translation-invariant solutions of the field equations
% BANKS-SOURCE: pdf=205 print=195 kind=equation
\[
  \frac{\delta\Gamma}{\delta\phi(x)}=0
\]
that violate symmetries of the classical action. To do this, we need only
calculate the \emph{effective potential}, which is defined by\footnote[20]{We
work in Euclidean space. In Lorentzian signature this equation needs a minus
sign.}
% BANKS-SOURCE: pdf=205 print=195 kind=equation
\[
  \Gamma[\phi_c]=\frac{1}{g^2}\int\dd^D x\,V_{\mathrm{eff}}(\phi_c).
\]
$g^2$ is the loop-counting parameter. The derivatives of $V_{\mathrm{eff}}$
with respect to $\phi_c$ are 1PI Green functions evaluated at zero momentum.

To develop an efficient graphical method to calculate $V_{\mathrm{eff}}$, we
begin with its Legendre transform, the connected generating functional in a
constant classical source $J$, divided by the volume of space-time. $w(J)$ is
computed as the logarithm of a path integral, in the usual way, except that we
divide the answer by the volume, and multiply by $g^2$. At tree level, $w(J)$
is just the Legendre transform of the potential in the classical action. More
generally, we can think of it as the energy density of the ground state in the

% BANKS-SOURCE: pdf=206 print=196 kind=source-page
presence of the constant classical source. There are two interesting
interpretations of $V_{\mathrm{eff}}$, which follow from these definitions.

First of all, it is the answer to the following variational question: ``What is
the minimal value of the expectation value of the Hamiltonian of the $J=0$
theory, in states constrained to satisfy
% BANKS-SOURCE: pdf=206 print=196 kind=display
\[
  \langle\psi|\phi|\psi\rangle=\phi_c?''
\]
Indeed, the way to solve this constrained variational problem is to introduce
$J$ as a Lagrange multiplier, calculate the minimal energy as a function of $J$
(which is $w(J)/V_{\mathrm{space}}$) and then impose the constraint
% BANKS-SOURCE: pdf=206 print=196 kind=display
\[
  w'(J)=\phi_c.
\]
The minimal expectation value of $H$ is then obtained by subtracting the
source term from $w$ and then expressing the answer as a function of $\phi_c$.
This is just the Legendre transform.

Consider instead the path integral
% BANKS-SOURCE: pdf=206 print=196 kind=display
\[
  \int[d\phi]\ee^{\ii S}\delta\!\left(\int\phi
  -V_{\mathrm{space-time}}\phi_c\right).
\]
Write the delta function as
% BANKS-SOURCE: pdf=206 print=196 kind=display
\[
  \int\dd J\,\ee^{\ii(\int\phi J-V_{\mathrm{space-time}}\phi_cJ)},
\]
and perform the path integral to obtain
% BANKS-SOURCE: pdf=206 print=196 kind=display
\[
  \int\dd J\,\ee^{\ii V_{\mathrm{space-time}}(w(J)-J\phi_c)}.
\]
Now do the integral over $J$ by stationary phase in the large-volume limit. The
answer is
% BANKS-SOURCE: pdf=206 print=196 kind=display
\[
  \ee^{-\ii V_{\mathrm{space-time}}V_{\mathrm{eff}}(\phi_c)}.
\]
In other words, in the large-volume limit, the effective Wilsonian action for
the zero-momentum mode of the field is just the effective potential.

In order to compute $V_{\mathrm{eff}}$ perturbatively we imagine that we have
found the value of $J$ corresponding to the expectation value of $\phi$ equal
to $\phi_c$. Then the effective potential is just $w(J)-J\phi_c$. Now shift
the field in the path integral defining $w(J)$ by $\phi_c$: $\phi=\phi_c+\Delta$.
By definition, the field $\Delta$ has zero expectation value. The perturbative
way to determine $J=\sum g^nJ_n$ is to insist that it is chosen so that, in
each order, the linear term from the explicit $J\Delta$ term in the action
cancels out the ``tadpole'' graphs which give the one-point function of
$\Delta$. This then implies that all one-particle-reducible (1PR) vacuum graphs vanish. We
thus obtain the following prescription \cite{banks-ref-125}.

Define the action
$S(\phi_c+\Delta)-V_{\mathrm{space-time}}\phi_c(\partial V/\partial\phi)(\phi_c)$.
Compute the vacuum path integral for this action, ignoring all 1PR diagrams.
This path integral is equal to
% BANKS-SOURCE: pdf=206 print=196 kind=display
\[
  \ee^{-V_{\mathrm{space-time}}V_{\mathrm{eff}}(\phi_c)}.
\]

% BANKS-SOURCE: pdf=207 print=197 kind=source-page
The one-loop term in this expression needs some care, since it is unusual to
look at a Feynman diagram with no vertices. Let's consider an $O(n)$-invariant
Euclidean Lagrangian
% BANKS-SOURCE: pdf=207 print=197 kind=display
\[
  \mathcal{L}=\frac12(\partial_\mu\phi^i)^2+V(\phi^i\phi^i).
\]
According to our prescription, the one-loop correction to $V_{\mathrm{eff}}$
is
% BANKS-SOURCE: pdf=207 print=197 kind=display
\[
\begin{aligned}
  V_{\mathrm{eff}}^{(1)}&=-\ln\left(\int[d\Delta]\,
  \ee^{-\frac12(\Delta^i(-\nabla^2)\Delta^i+(M^2)_{ij}\Delta^i\Delta^j)}\right)\\
  &=\frac12\ln\det(-\nabla^2+M^2).
\end{aligned}
\]
The matrix $M^2$ is defined by
$M_{ij}^2=(\partial^2V/\partial\phi^i\partial\phi^j)(\phi_c)$.

To find the regularized form of this answer, let us divide the functional
integral by the same expression with $M^2\to M_0^2\delta_{ij}$, where $M_0$
is a large mass that we think of as order the cut-off. This subtracts a
$\phi_c$-independent constant from $V_{\mathrm{eff}}$ and leads to a simpler
formula at intermediate stages of the calculation. It has no relevance to any
physical quantity. For any Hermitian matrix, $\operatorname{tr}\ln A=\ln\det A$,
so we can write
% BANKS-SOURCE: pdf=207 print=197 kind=display
\[
  V_{\mathrm{eff}}^{(1)}=-\frac12\operatorname{Tr}\left(\int\frac{\dd s}{s}
  \left[\ee^{-s(-\nabla^2+M^2)}-\ee^{-s(-\nabla^2+M_0^2)}\right]\right).
\]
The identity which leads to this parametric formula can be verified eigenvalue
by eigenvalue for the two operators. Note that it looks just like the
Schwinger proper-time version of a one-loop vacuum diagram with no vertices,
except that the proper-time integral has a factor of $1/s$ in it that we might
not have guessed from ordinary Feynman rules. The trace
$\operatorname{Tr}O=\int\dd^{d_{\mathrm{reg}}}x\langle x|\operatorname{tr}O|x\rangle$ is over the
tensor product of function space and $O(n)$-index space; $\operatorname{tr}$
refers to the index trace alone. The operator $-\nabla^2+M^2$ is diagonal in
Fourier space. Thus we obtain
% BANKS-SOURCE: pdf=207 print=197 kind=display
\[
\begin{aligned}
  V_{\mathrm{eff}}^{(1)}={}&-\frac12V_{\mathrm{space-time}}\operatorname{tr}
  \left(\int\frac{\dd^{d_{\mathrm{reg}}}p}{(2\pi)^{d_{\mathrm{reg}}}}\frac{\dd s}{s}
  \left[\ee^{-s(p^2+M^2)}-\ee^{-s(p^2+M_0^2)}\right]\right)\\
  ={}&-\frac{1}{2(2\sqrt\pi)^{d_{\mathrm{reg}}}}
  \operatorname{tr}M^{d_{\mathrm{reg}}}
  \Gamma\!\left(-\frac {d_{\mathrm{reg}}}{2}\right).
\end{aligned}
\]
Now
% BANKS-SOURCE: pdf=207 print=197 kind=display
\[
  \Gamma\!\left(-\frac {d_{\mathrm{reg}}}{2}\right)=
  \frac{1}{(2-d_{\mathrm{reg}}/2)(1-d_{\mathrm{reg}}/2)(-d_{\mathrm{reg}}/2)}
  \Gamma\!\left(3-\frac {d_{\mathrm{reg}}}{2}\right),
\]
so the one-loop effective potential has a pole at $d_{\mathrm{reg}}=4$. If the tree-level
potential is quartic in the fields,
$V=(m_0^2/2)\phi^2+\lambda_0(\phi^2)^2$, then the residue of this pole
proportional to $\operatorname{tr}M^4$ is a mixture of quadratic and quartic
terms. Furthermore, it is an $O(n)$-invariant function of the classical fields
$\phi$. This means that the pole can be removed by redefining $m_0$ and
$\lambda_0$. The required renormalization (in the MS scheme) is identical to
that which we would have to do in order to eliminate the divergent part of any
one-loop Green function.

% BANKS-SOURCE: pdf=208 print=198 kind=source-page
This calculation also shows us that theories with spontaneously broken
symmetries are renormalized by the same manipulations as those which renormalize
the corresponding symmetries in the symmetric phase. Indeed, since we are
calculating the quantum action, we never have to make a choice of the vacuum we
expand around. The renormalized quantum action is an $O(n)$-invariant
functional of $\phi$. The renormalized parameters (which, as low-energy
observers, we are free to choose at will) are such that the minimum of the
effective potential is asymmetric.\footnote[21]{Strictly speaking, the exact
effective potential is a convex function of the field. If spontaneous symmetry
breaking occurs, it is in fact constant over the region $|\phi|<\phi_0$,
indicating that the true vacua are rotations of the vector $\phi_0(1,0,\ldots,0)$.
Any point in the indicated region can be the VEV of the field in appropriate
superpositions of these rotated vacua. Only the vacua with fixed rotation
angles satisfy the cluster-decomposition principle. One never sees the
convexity of the effective potential in simple approximations.} This decoupling
between renormalization and spontaneous symmetry breaking should have been
expected. We described spontaneous breaking in terms of inequivalent
representations of the same field algebra, which were inequivalent precisely
because of their different large-volume behaviors. Renormalization has to do
with the local structure of the theory.

% BANKS-SOURCE: pdf=208 print=198 kind=subsection id=9.14.5
\subsubsection{The Coleman--Weinberg potential}

The computation of the effective potential in a general renormalizable field
theory follows the pattern set by the $O(n)$ model. One simply computes the
vacuum energy as a function of particle masses, where the masses are those
induced by a scalar field VEV $v$. The resulting renormalized potential in $D=4$ reads
% BANKS-SOURCE: pdf=208 print=198 kind=equation
\[
  V(v)=\frac{1}{64\pi^2}\sum_i(2S+1)(-1)^{2S}M_i^4(v)
  \ln\!\left(M_i(v)^2/\mu^2\right),
\]
where $\mu$ is the renormalization scale. This is evaluated in terms of
renormalized couplings and satisfies a renormalization-group equation. The
couplings (and Lagrangian mass parameters) come into the computation of the
mass eigenvalues $M_i(v)$. One generally has to diagonalize a complicated mass
matrix in order actually to evaluate this formula. Note that the minus sign for
fermions comes about because Grassmann functional integrals give determinants
rather than inverse determinants. More intuitively it is a consequence of the
fact that ``vacuum fermions have negative energy.''

It is interesting to note that this expression is analytic in $v$, except at
places where one or more mass eigenvalues go to zero. We have obtained an
effective action for the zero-momentum modes of the fields, by integrating out
the higher modes. However, if a mass is zero, there is no clear separation.
This shows up as IR divergences in loop diagrams. The connection between
non-analyticity of the effective potential and massless particles was the key
to Wilson's analysis of critical phenomena in condensed-matter physics
\cite{banks-ref-126}. It is also of crucial importance in particle physics. One
typical example is the Linde--Weinberg lower bound on the mass of the Higgs
particle in the standard model \cite{banks-ref-127,banks-ref-128}. If the Higgs
is much lighter than the W and Z bosons, they can be integrated out and
contribute only to the Higgs effective action. At tree level,

% BANKS-SOURCE: pdf=209 print=199 kind=source-page
the ratio of masses is essentially the ratio between $\sqrt\lambda$ and
$g_{1,2}$. However, no matter how small the Higgs quartic coupling is, compared
with the gauge couplings, the gauge bosons become massless at $|H|=0$. Their
non-analytic contribution to the Coleman--Weinberg potential dominates the
tree-level potential in this region, and produces a symmetric minimum if
$\lambda$ is very small. The condition that the Universe actually lives in the
spontaneously broken vacuum state then puts a bound on $m_H/m_W$.

% BANKS-SOURCE: pdf=209 print=199 kind=subsection id=9.14.6
\subsubsection{Renormalization of gauge theories in the Higgs phase}

The interaction of renormalization and the Higgs phenomenon is more subtle than
that between renormalization and spontaneously broken global symmetries. In the
unitary gauge, the gauge-boson propagator does not fall off at high energies and
non-abelian Higgs models are not renormalizable in unitary gauge. However, the
$S$-matrix in this gauge is equal to that in the general $R_\kappa$ gauge, which
is renormalizable.

This peculiar situation can be understood better by noting that the elementary
fields in unitary gauge are equal to gauge-invariant operators
% BANKS-SOURCE: pdf=209 print=199 kind=equation
\[
  B_\mu=\Omega^\dagger D_\mu(A)\Omega,
\]
etc., where $\Omega$ is defined by
% BANKS-SOURCE: pdf=209 print=199 kind=equation
\[
  H(x)=\Omega(x)(v+\Delta(x)).
\]
$v$ is the VEV of $H$, and $\Delta$ (the vector of physical Higgs fields)
satisfies $v^TT_S^a\Delta=0$. These gauge-invariant operators can be defined
in any gauge, but are equal to the elementary fields of finite mass in the
$\kappa\to\infty$ limit.

Thus, in $R_\kappa$ gauge for general $\kappa$, the elementary fields of the
unitary gauge are defined as non-polynomial functions of the elementary fields
of the $R_\kappa$ gauge. In perturbation theory, we are always working around
the Gaussian fixed point, where the operators of fixed dimension are defined as
Wick monomials. Higher-order monomials require more subtractions. Thus, it is
not surprising that the Green functions of unitary gauge fields require an
infinite number of subtractions, which is characteristic of a non-renormalizable
theory.

On the other hand, we have argued that the scattering matrix can be computed
from projected Green functions of the elementary fields, by use of the LSZ
formula. These projected operators, smeared with normalizable wave packets at
infinity, are BRST-invariant. Thus, the $S$-matrix and the Green functions of
those gauge-invariant operators that can be defined as polynomials in the
$R_\kappa$ gauge fields behave as we expect a renormalizable theory to behave.

This remark also sheds light on the situation we encountered when computing in
massive QED. There we found an apparent quadratic UV divergence, inversely
proportional to the square of the gauge-boson mass, in the wave-function
renormalization of the fermion field. When we used the gauge-invariant DR
procedure, this divergence disappeared and was replaced by a logarithmic
divergence with no singularity at $\mu=0$. We could have done the calculation
using the Stueckelberg formalism in an $R_\kappa$ gauge,

% BANKS-SOURCE: pdf=210 print=200 kind=source-page
where the gauge-invariant regulator would have been natural. We had to use the
gauge-invariant regulator in the unitary gauge in order to get a result
consistent with the Higgs mechanism. Note that the unitary gauge abelian Higgs
model is renormalizable, while this is not true for non-abelian theories. The
difference has to do with the fact that the effective theory of an abelian NGB
is free, whereas that of a non-abelian NGB is a non-renormalizable interacting
theory.

% BANKS-SOURCE: pdf=210 print=200 kind=subsection id=9.14.7
\subsubsection{Effective field theory and NGBs}

Theories of quantum gravity do not have continuous global internal symmetries.
This is a theorem in perturbative string theory \cite{banks-ref-129}, and
follows more generally from the violation of global conservation laws by
black-hole formation and evaporation. Thus, continuous global symmetries should
be viewed as accidental symmetries of low-energy effective field theories. As
such, we expect them to be broken by higher-dimension operators with scale no
higher than the reduced Planck scale $m_P=2\times10^{18}\,\mathrm{GeV}$. If $G$
is a spontaneously broken U(1) (subgroup of a) continuous accidental symmetry
group, and $D$ is the dimension of $O_D$, the lowest-dimension G-violating
operator allowed by exact symmetries of the underlying gravity theory, then we
expect a potential for the approximate NGB field $b$ associated with $G$, of
the form
% BANKS-SOURCE: pdf=210 print=200 kind=equation
\[
  V=\frac{\langle O_D\rangle}{m_P^{D-4}}U(b/f).
\]
$f$ is the PNGB decay constant. When $G$ is something like a global chiral
symmetry of a strongly coupled gauge theory, we expect $\langle O_D\rangle
\sim f^D$. The approximate NGB mass is then of order
% BANKS-SOURCE: pdf=210 print=200 kind=display
\[
  m_b\sim\left(\frac{f}{m_P}\right)^{(D-4)/2}f.
\]

A context in which these effects are important is the axion solution of the
strong CP problem. The axion is a pseudo-NGB, whose potential is supposed to
come dominantly from a coupling $(a/f_a)G\widetilde G$, to the gluons of QCD
\cite{banks-ref-130}. In most models where the axion arises from spontaneous
breaking of a new strongly coupled gauge theory, the mass arising from
``Planck slop'' is larger than that coming from QCD, and the axion model does
not work.

In general, even if we neglect symmetry-breaking terms, the low-energy effective
theory for NGBs of a non-abelian group is not renormalizable. The interactions
are all irrelevant operators. With our current understanding of renormalization,
this should not bother us. The constant which sets the scale for all these
irrelevant interactions is $f$, the NGB decay constant. The full effective
theory contains all possible irrelevant operators consistent with the
symmetries. In addition one is instructed to calculate with a cut-off $4\pi f$
(in four dimensions). To simplify matters one chooses the cut-off procedure to
preserve the symmetry which gave rise to the NGB in the first place.\footnote[22]{In
the first few paragraphs of this subsection we argued that all such symmetries
were approximate. The more precise meaning of the procedure we are now
outlining is that the irrelevant operators which break the symmetry are scaled
by $m_P$, or some other scale $\gg f$. As a consequence they contribute much
less than do the irrelevant terms we are discussing here, at energy scales
below $4\pi f$.}

% BANKS-SOURCE: pdf=211 print=201 kind=source-page
It is then easy to verify that, as long as NGB momenta are below $f$, the
higher-order terms in the effective Lagrangian make small corrections, in a
systematic power series in $p/f$, to amplitudes computed using the leading-order
effective action.

\BanksImplicitHook{I-CH09-023}

We can use this effective field theory to do loop calculations, and we must do
so in order to capture non-analytic terms in the amplitudes which arise from
loops of massless NGBs. All analytic corrections from the loops are simply
absorbed into a redefinition of the coefficients of higher-dimension operators
in the effective action. A systematic exposition of the rules of effective
field theory for NGBs may be found in Chapter 19 of \cite{banks-ref-131}.
